\problemname{Spacewalk}
Astronaut Gustav is serving on a space station consisting of $n$ modules connected in a circle so that module $1$ is connected to module $2$, $2$ to $3$, and so on (and module $n$ is connected to module $1$).
The distance between two adjacent modules is $1$. To create a sort of artificial gravity, the space station rotates at a constant speed around the center of the circle.

The station has been in space for quite a while and it is time to polish the outside of the windows.
The lot has fallen on Gustav to do this. There are $m$ windows numbered from $1$ to $m$, where window number $i$ is on module number $a_i$.
For some reason, the windows must be polished in exactly this order. The only entrance and exit to the station is at module $1$.

To travel between modules, there is a rocket-powered window elevator that runs along the outside of the station.
The window elevator can only travel between adjacent modules, and thus cannot take any shortcuts.
Gustav wants to choose a path from module $1$, around to all the windows, and back to module $1$.
Unfortunately, there are two problems: first, the elevator has limited fuel, so Gustav must choose a path that minimizes the distance he travels.
Furthermore, the elevator's movements affect the space station's rotation, so it must travel equally much clockwise as counterclockwise.

Find the minimum distance Gustav can travel so that he starts at module $1$, visits all windows in the correct order, returns to module $1$, and travels equally much counterclockwise as clockwise.

\section*{Input}
The first line contains two integers $n$ and $m$ ($3 \leq n \leq 10^5$, $1 \leq m \leq 10^5$), the number of modules and the number of windows.
The second line contains $m$ integers, the indices of the windows to be polished ($1 \leq a_i \leq n$).

\section*{Output}
Output one integer, the minimum possible distance.

\section*{Scoring}
Your solution will be tested on a set of test groups. To earn points for a group, you must pass all test cases in that group.

\noindent
\begin{tabular}{| l | l | p{12cm} |}
  \hline
  \textbf{Group} & \textbf{Points} & \textbf{Constraints} \\ \hline
  $1$    & $10$      & $3 \le n \le 10$, $1 \le m \le 10$ \\ \hline
  $2$    & $24$      & $3 \le n \le 100$, $1 \le m \le 100$ \\ \hline
  $3$    & $19$      & $1 \le m \le 1000$  \\ \hline
  $4$    & $47$      & No additional constraints. \\ \hline
\end{tabular}

\section*{Explanation of Sample 1}
The optimal path here is $1-2^*-3-4^*-3^*-4-5-6^*-5-4-3-2-1$ (the asterisk means a window is polished).
\section*{Explanation of Sample 2}
In this example, there is first a window on module $1$, which Gustav can polish without having to go anywhere. Then he goes to module 2,
polishes the three windows that are there, and goes back.
